\section{Preliminaries}\label{sec: prelim}

We are interested in polynomials over $\GF[2]$.
Most polynomials we use are of
degree $d$ but also polynomials of degree one, that
we call ``affine forms'' play a special role.
We are not interested in the polynomials as formal
polynomials, but rather as functions mapping
$\GF[2]^n$ to $\GF[2]$ and hence we freely
use that $x_i^2=x_i$ and thus any
monomial in our polynomials can be taken to be multi-linear.
We start with the following standard result which we,
for completeness, even prove.

\begin{theorem}\label{thm basic deg d}
Any multivariate polynomial $P$ of degree
$d$ that is nonzero takes the
value 1 for at least a fraction
$2^{-d}$ of the inputs.
\end{theorem}


\begin{proof}
The proof is by induction over $n$ and $d$,
with the base case of $d=1$ which is true
as each linear polynomial is unbiased.

For the induction step, suppose $P(x)=P_0(x)+x_1 P_1(x)$
and let us consider what happens for the two possible
values of $x_1$.
If both $P_0$ and $P_0+P_1$ are non-zero we are
done by induction.  If not, as $P_1$ is of degree
at most $d-1$, the polynomial of $P_0$ and $P_0+P_1$ 
that is non-zero is of degree at most $d-1$.
Hence this polynomial takes the value 1 for at least a fraction $2^{1-d}$
of \emph{its} inputs.  As the set of inputs of this
polynomial constitutes half of the inputs of $P$,
the result follows also in this case.
\end{proof}

It is not difficult to see that this result is
tight by considering $P(x)=\prod_{i=1}^dx_i$,
or, more generally, products of $d$ linearly independent
affine forms.  
It is important for us
that these are the only cases of tightness.
This follows from a characterization by Kasami
and Tokura~\cite{kt70} of all polynomials
that are non-zero for at most a fraction
$2^{1-d}$ of the inputs.  
A consequence of their characterization is
the following theorem. 

\begin{theorem}\label{basic theorem}
Let $P$ be a %
degree-$d$
polynomial
over $\GF[2]$ which factors
as
$$P(x)= Q(X) \prod_{i=1}^r A_i (x) 
$$
where $A_i$ are linearly independent
affine forms and $Q$ does not
contain any affine factor.
Then the fraction of points on which $P(x)=1$
is at least
$$
2^{-r} (2^{1-(d-r)}-2^{1-2(d-r)})\,,$$
if $d\not =r$ and $2^{-d}$ if $d=r$.
\end{theorem}

Given that we do not need their full characterization
and the proof of the part that we need is shorter
than the proof of the full characterization, we
give the proof of \expref{Theorem}{basic theorem} in
an appendix.
